\section*{第 \ref{ch:g5:simple-machines} 章 --- 简单机械：杠杆、滑轮、斜面}

\begin{solution}{exo:g5:simple-machines:1}
杠杆（撬棍、跷跷板、剪刀）、滑轮（旗杆、水井、
起重机）、斜面（轮椅坡道、山路、
装货坡道）。
\end{solution}

\begin{solution}{exo:g5:simple-machines:2}
定滑轮只改变拉力的方向（向下拉，
重物向上升）。动滑轮改变的是力的大小：重物由两段绳子悬挂，
所以手只需承受一半的重力。
\end{solution}

\begin{solution}{exo:g5:simple-machines:3}
向下拉时可以借助自身的重力：你可以倚着绳子，甚至
整个人挂在绳子上，让地球的引力帮助肌肉用力。
向上提时，只能靠手臂单独用力。
\end{solution}

\begin{solution}{exo:g5:simple-machines:4}
拉力只需重力的一半。要把重物提升 \qty{3}{m}，两段绳子都必须
各缩短 \qty{3}{m}：你的手需要拉过 $2 \times 3 = \qty{6}{m}$ 的
绳子。
\end{solution}

\begin{solution}{exo:g5:simple-machines:5}
陡坡所需的拉力更大。
缓坡省力，代价是路程变长：同样的高度，
要沿着长得多的路才能达到，这就是黄金法则的交换。
\end{solution}

\begin{solution}{exo:g5:simple-machines:6}
若直接沿斜坡向上拉，所需的拉力大到连任何发动机或
骡队都无法承受；之字形道路把一段又长又缓的斜面
铺展在山腰上——用力小，路程长。整条路就是一个
被拉长了的斜面。
\end{solution}

\begin{solution}{exo:g5:simple-machines:7}
任何机械都不会“白送”：省下的力，
必须以距离严格偿还。力为十分之一，所付出的路程就是十倍，
（或绳子）。
\end{solution}

\begin{solution}{exo:g5:simple-machines:8}
螺纹就是一圈绕在杆上的又长又缓的斜面。螺丝刀的许多次
轻松转动就是那段长路；每转一下螺钉的微小下沉，
就是所获得的小高度。
\end{solution}

\begin{solution}{exo:g5:simple-machines:9}
拉力只需重力的四分之一。四段绳子各要缩短
\qty{2}{m}：共 $4 \times 2 = \qty{8}{m}$ 的绳子
要穿过工人的双手。
\end{solution}

\begin{solution}{exo:g5:simple-machines:10}
新斜面爬升同样的 \qty{1}{m}，但路程只有一半，因此坡度加倍，
推力也大致加倍。黄金法则
当场得到体现：路程减半，力加倍——箱子
的爬升所做的功无论哪种方式都相同，搬运工
不过是用许多轻松的小步，换成了几步沉重的大步。
\end{solution}

\begin{solution}{exo:g5:simple-machines:11}
任何安排都无法兑现这一承诺。把重物升高一米需要固定的
能量；即使只用一半的力，卖家
也必须拉过\emph{两}米绳子——只用一半的力沿一米路程拉，
只能交回一半的账单。广告许诺的是
半价的“有偿”提升，这正是黄金法则与
那些失败的永动机所禁止的“免费午餐”。
\end{solution}
